% Solutions to the exercises of Chapter 17. Numbers: code/solutions/sol17.py
\documentclass[../solutions_manual.tex]{subfiles}
\begin{document}
\begingroup\setlength{\emergencystretch}{3em}% local to this chapter
\ifSubfilesClassLoaded{\setcounter{chapter}{16}}{}% NN-1 when built standalone
\setcounter{chapter}{16}% NN-1, so this chapter is numbered NN
\chapter{Domestic Structural Diagnostics}\label{sol17}

Every number below is printed by \texttt{python3} \nolinkurl{code/solutions/sol17.py}. Input-output results use INE's supply-use tables \citep{ine2026sut} as processed in \nolinkurl{data/processed/bolivia_Z} and \nolinkurl{bolivia_io_*.csv}; comparators use OECD ICIO 2023 \citep{oecd2023icio}; employment uses ILOSTAT \citep{ilo2026ilostat}. All were retrieved on 2026-11-05. Household-survey microdata (EH) and audience data were \textbf{not fetched} (\nolinkurl{data/README.md}, \S9). Exercises that need them get the method only.

%% ---------------------------------------------------------------------
\exercisesolution{17.1}{Compute Bolivia's $\Gamma$}

\textbf{Method.}
\begin{enumerate}
\item INE's 2019 supply-use tables (base 2017) are converted to a domestic industry$\times$industry table at basic prices: import proportionality, product-sales-structure model D (\nolinkurl{data/README.md}, \S2).
\item Aggregate to the book's 12 sectors, rows = sellers.
\item $a_{ss'} = Z_{s's}/X_s$ (Appendix~17A); $W^P = D^{-1}A$, so row $s$ is $Z_{\cdot s}/\sum_{s''}Z_{s''s}$.
\item Domestic services has no intermediate purchases or sales, so its row is undefined and it is dropped; $N = 11$.
\end{enumerate}

\textbf{Results, input shares (book12).} The eigenvalues by modulus are $1$; $0.5326$; $0.3127$; $0.2233$; $-0.1951 \pm 0.0986i$; $0.1003 \pm 0.1035i$; $-0.0192$; $0.0171$; $-0.0137$. Hence $\lambda_2 = 0.5326$ (real, simple) and $\Gamma = 0.4674$. The Perron vector is
\begin{center}
\small
\begin{tabular}{lc|lc|lc}
\toprule
Hydrocarbons & 0.039 & Construction & 0.029 & Agriculture & 0.091 \\
Mining & 0.037 & Commerce & 0.095 & Other services & 0.270 \\
Public administration & 0.005 & Transport & 0.092 & Personal services & 0.021 \\
Manufacturing & 0.159 & Financial services & 0.162 & & \\
\bottomrule
\end{tabular}
\end{center}
so $G(\bpi) = 0.4446$ and $\Neff = 6.48$. Checks: $\sum\pi = 1$; all values agree with Examples~\ref{ch17:ex:bolivia_gamma} and~\ref{ch17:ex:bolivia_pi} and with \nolinkurl{bolivia_io_summary.csv} to $10^{-9}$.

\textbf{Normalizations and aggregations} (2019):
\begin{center}
\small\setlength{\tabcolsep}{3pt}
\begin{tabular}{llcccl}
\toprule
Aggregation & Normalization & $\Gamma$ & $G(\bpi)$ & $\Neff$ & $\arg\max\pi$ \\
\midrule
book12 & input shares & 0.467 & 0.445 & 6.48 & Other services (0.270) \\
book12 & input, zero diagonal & 0.614 & 0.363 & 7.79 & Other services (0.190) \\
book12 & sales shares & 0.694 & 0.319 & 7.88 & Manufacturing (0.241) \\
book12 & sales, zero diagonal & 0.653 & 0.268 & 8.55 & Manufacturing (0.221) \\
book12\_isic & input shares & 0.475 & 0.435 & 6.67 & Other services (0.258) \\
8 sectors (here) & input shares & 0.589 & 0.517 & & Finance \& other serv.\ (0.416) \\
8 sectors (here) & sales shares & 0.676 & 0.237 & & \\
detail39 & input shares & 0.410 & 0.670 & & (CSV, $N = 38$) \\
detail39 & sales shares & 0.331 & 0.428 & & (CSV) \\
\bottomrule
\end{tabular}
\end{center}
The 8-sector grouping merges hydrocarbons with mining, commerce with transport, finance with other services, and personal with domestic services.

\textbf{How much do they change?} Across book12 and detail39, input or sales shares, with or without diagonal, $\Gamma$ ranges over $0.331$--$0.694$ (the book's $0.33$--$0.69$) and $G(\bpi)$ over $0.268$--$0.670$. The spread across normalizations, about $0.36$ in $\Gamma$, is wider than the spread across the 23 economies of Exercise~16.6 ($0.333$--$0.659$). Three regularities:
\begin{itemize}
\item Zeroing the diagonal raises $\Gamma$, because large self-purchases (finance $0.567$, other services $0.511$) are the slow modes.
\item Sales shares raise $\Gamma$ and lower $G(\bpi)$: the hub becomes manufacturing as a buyer rather than other services as a supplier.
\item Aggregating to 8 sectors raises $\Gamma$, but merging finance with other services creates a single large hub, which raises $G(\bpi)$.
\end{itemize}

\emph{Model answer.} Report $\Gamma$ with its normalization and aggregation, and compare countries only like-for-like (book12\_isic, domestic, input shares). Under any variant, no eigenvalue of Bolivia's $W^P$ is close to $1$, so there are no weakly coupled blocks at this resolution.

%% ---------------------------------------------------------------------
\exercisesolution{17.2}{Between-block variance}

\emph{Data status.} EH 2019 microdata were \textbf{not fetched} (INE ANDA portal; row 41 of the status table in \nolinkurl{data/README.md}).

\solpart{a--b} \textbf{Method.}
\begin{enumerate}
\item From the EH 2019 person file, keep employed adults with positive labour income. Use the survey weights throughout; monthly labour income, or per-capita household income for the household version.
\item Assign blocks: (1) workers in ISIC B (mining and hydrocarbons) or O (public administration); (2) formal urban, i.e.\ urban workers contributing to social security or employed by firms with 5 or more workers; (3) everyone else (informal or rural).
\item With weights $w_i$, compute $\bar y = \sum_i w_iy_i/\sum_i w_i$, block shares $N_b/N = \sum_{i\in b}w_i/\sum_i w_i$, block means $\bar y_b$, $V_B$ by~\eqref{ch17:eq:between} and $V_T = \sum_i w_i(y_i - \bar y)^2/\sum_i w_i$. Then $\sigma_B = V_B/V_T$.
\item Bootstrap $\sigma_B$ by resampling primary sampling units.
\end{enumerate}
With the chapter's illustrative inputs (Example~\ref{ch17:ex:bolivia_between}): $\bar y = 4{,}326$, $V_B = 19.98$M and $\sigma_B = 0.624$ for $V_T = 32$M.

\solpart{c} \textbf{Two and four blocks.} Merging blocks 1 and 2 of the illustration into ``formal'' gives a formal mean of $7{,}674.4$, $V_B = 8.46$M and $\sigma_B = 0.264$: the coarser partition loses $11.52$M of between-block variance. That loss is exactly the variance of the two fine means inside the merged block, as the proof below shows. A four-block partition (gas, formal manufacturing, urban services, rural) can only raise $\sigma_B$ relative to any three-block partition it refines. \emph{Expected with real data:} $\sigma_B$ rises with the number of blocks, at a decreasing rate.

\solpart{d} \textbf{Refinement never lowers $\sigma_B$.} Let $\mathcal{P}'$ refine $\mathcal{P}$: every block $B$ of $\mathcal{P}$ is a union of blocks $b$ of $\mathcal{P}'$. With $n_b$ the size of fine block $b$ and $n_B = \sum_{b\subseteq B}n_b$, the coarse mean is $\bar y_B = \sum_{b\subseteq B}(n_b/n_B)\bar y_b$. Then
\[
\textstyle V_B(\mathcal{P}') = \sum_B\sum_{b\subseteq B}\frac{n_b}{N}(\bar y_b - \bar y)^2 = V_B(\mathcal{P}) + \sum_B\frac{n_B}{N}\sum_{b\subseteq B}\frac{n_b}{n_B}(\bar y_b - \bar y_B)^2 \;\ge\; V_B(\mathcal{P}).
\]
The middle equality writes $\bar y_b - \bar y = (\bar y_b - \bar y_B) + (\bar y_B - \bar y)$ and notes that the cross term vanishes within each $B$ (the law of total variance applied to block means). $V_T$ does not depend on the partition, so $\sigma_B$ weakly rises under refinement, with equality iff all fine means inside every coarse block are equal. The highest $\sigma_B$ is reached by the finest partition: with singletons, $V_W = 0$ and $\sigma_B = 1$. The script checks both facts on $500$ random samples.

\emph{Implications.} $\sigma_B$ is comparable across countries only for partitions of the same number and type of blocks. There is no ``natural'' partition in the data unless one is fixed in advance on substantive grounds, or read off $W^P$ by spectral clustering (Appendix~17B). A partition chosen by maximizing $\sigma_B$ is meaningless: the maximum is always the singleton partition.

%% ---------------------------------------------------------------------
\exercisesolution{17.3}{Power-employment ratio}

\solpart{a} $\bpi$ is that of Exercise~17.1 (input shares, book12, 2019).

\solpart{b} \textbf{Ratios.} Employment is from ILOSTAT 2019 (ISIC Rev.~4 sections, mapped as in Example~\ref{ch17:ex:bolivia_pi}). VA per worker is INE value added divided by ILOSTAT employment, in thousand Bs per worker.
\begin{center}
\small
\begin{tabular}{lcccc}
\toprule
Sector & Emp.\ share & $\pi_s$ & $\pi_s/\text{share}_s$ & VA per worker \\
\midrule
Hydrocarbons and mining & 0.013 & 0.077 & 5.74 & 464.6 \\
Public administration & 0.026 & 0.005 & 0.19 & 313.0 \\
Manufacturing & 0.112 & 0.159 & 1.42 & 60.4 \\
Construction & 0.079 & 0.029 & 0.37 & 54.6 \\
Commerce & 0.171 & 0.095 & 0.55 & 29.3 \\
Transport & 0.071 & 0.092 & 1.30 & 39.1 \\
Financial services & 0.007 & 0.162 & 23.36 & 378.1 \\
Agriculture & 0.274 & 0.091 & 0.33 & 16.8 \\
Other services & 0.121 & 0.270 & 2.23 & 81.7 \\
Personal services & 0.107 & 0.021 & 0.20 & 41.3 \\
Domestic services & 0.019 & --- & --- & 19.0 \\
\bottomrule
\end{tabular}
\end{center}

\solpart{c} \textbf{Ranking and capital intensity.} The highest ratio is financial services ($23.4$); the lowest is public administration ($0.19$), then personal services ($0.20$) and agriculture ($0.33$). Regressing $\log(\pi/\text{share})$ on $\log(\text{VA per worker})$ over the 10 sectors gives a slope of $0.770$ and $R^2 = 0.328$ (Spearman $0.515$). So value added per worker, a proxy for capital intensity and labour productivity, explains about a third of the variation.

The rest has a clean decomposition:
\[
\frac{\pi_s}{\text{emp}_s} = \frac{\pi_s}{\text{VA}_s/\text{VA}}\cdot\frac{\text{VA}_s/\text{emp}_s}{\text{VA}/\text{emp}}.
\]
The first factor, power per unit of value added, is $3.41$ for finance, $1.83$ for transport and $0.03$ for public administration. Public administration has the third-highest VA per worker but sells almost nothing to other sectors, so its residual is $-2.72$ log points. Finance has both a high VA per worker and a high power per unit of VA (residual $+1.93$). As the intuition box after Example~\ref{ch17:ex:bolivia_pi} says, the ratio mixes supplier centrality with productivity; it is not a measure of individual privilege.

\solpart{d} \textbf{Moving from agriculture to manufacturing.} Let $\varphi := \Phi/\Omega \in [0,1]$ and $N = 11$ sectors. Sector $s$ receives $x_s = \Omega[1/N + (\pi_s - 1/N)\varphi]$ (Proposition~\ref{ch07:prop:allocation-rule}), shared equally by its $\text{emp}_s\cdot E$ workers. Relative to the national mean per worker $\Omega/E$, a worker in $s$ gets $[1/N + (\pi_s - 1/N)\varphi]/\text{emp}_s$:
\begin{center}
\begin{tabular}{lcccc}
\toprule
$\varphi$ & 0 & 0.25 & 0.5 & 1 \\
\midrule
Agriculture & 0.332 & 0.332 & 0.332 & 0.332 \\
Manufacturing & 0.810 & 0.962 & 1.113 & 1.415 \\
Ratio & 2.44 & 2.90 & 3.35 & 4.26 \\
\bottomrule
\end{tabular}
\end{center}
At $\varphi = 1$ the relative allocation is exactly the power-employment ratio, so she gains a factor $1.42/0.33 = 4.26$. Agriculture's row does not move with $\varphi$ because $\pi_{\text{agr}} = 0.091 \approx 1/11$.

\emph{What the calculation assumes.}
\begin{enumerate}
\item Sectoral $\bpi(W^P)$ is a bargaining weight (Assumption~7.A). That is a modelling assumption, and this layer is the least natural one for it.
\item Sectors are the bargaining units and divide their share equally among workers.
\item One worker's move changes neither $W^P$ nor $\bpi$ (it is marginal).
\item Disagreement payoffs are equal across sectors.
\end{enumerate}
With a different $N$ (e.g.\ 39 sectors) the numbers change. \emph{Common mistake:} reading the factor $4.26$ as a wage premium; it is a division rule, not an estimate.

%% ---------------------------------------------------------------------
\exercisesolution{17.4}{Information democracy}

\emph{Data status.} Audience data were \textbf{not fetched} (Kantar IBOPE is paywalled; \nolinkurl{data/README.md}, \S9). Method only, plus a formal result on node sets.

\solpart{a--b} \textbf{Method.}
\begin{enumerate}
\item List outlets (TV channels, radio networks, newspapers, major digital platforms) with audience measures from a ratings service or a representative survey (``where did you get news yesterday?'').
\item Convert each outlet's reach to shares of total attention (minutes, or mentions as main source) that sum to $1$.
\item Group outlets into ecosystems for the four-node comparison.
\end{enumerate}
Using audience shares $\mathbf{s}$ as $\bpi^I$ is exact only if $W^I = \ones\mathbf{s}'$, i.e.\ everyone attends to the same mix, which has $\Gamma(W^I) = 1$ (Example~\ref{ch17:ex:bolivia_neff}). With survey data on \emph{who} watches what, build $W^I$ by groups and take its Perron vector instead.

\solpart{c} $\Neff = 1/\sum_j s_j^2$. For the chapter's four ecosystems, $\Neff = 3.381$.

\solpart{d} \textbf{Node sets.} \emph{Claim:} refining the node set never lowers $\Neff$. If node $e$ splits into outlets $o \in e$ with shares $s_o$ summing to $s_e$, then $\sum_{o\in e}s_o^2 \le (\sum_{o\in e}s_o)^2 = s_e^2$, since the cross terms are nonnegative. Summing over $e$, $\sum_o s_o^2 \le \sum_e s_e^2$, with equality iff no node is actually split. Examples:
\begin{itemize}
\item an equal split of every ecosystem into $k$ outlets multiplies $\Neff$ by $k$: $6.761$ for $k = 2$ and $16.903$ for $k = 5$;
\item an unequal split (0.8/0.2 for state media, 0.5/0.5 for private, none for community, 0.6/0.4 for social) gives $5.666$.
\end{itemize}
So ``Bolivia $3.4$ vs.\ the US $15$--$20$'' compares different node sets and says nothing. Compare on the same four-ecosystem partition, or report $\Neff/N$ on comparable outlet-level sets. \emph{Expected:} on four ecosystems most countries lie between 2 and 4; on outlet sets the numbers scale with the number of outlets measured.

%% ---------------------------------------------------------------------
\exercisesolution{17.5}{Bolivia's structural trajectory}

\solpart{a} \textbf{Every available year} (book12, domestic, input shares; from \nolinkurl{bolivia_io_summary.csv}, reproduced by the script from the flow matrices for the years in Example~\ref{ch17:ex:trajectory}):
\begin{center}
\scriptsize
\begin{tabular}{lcccc|lcccc}
\toprule
Year & $N$ & $\Gamma$ & $G(\bpi)$ & $\pi_{\text{hyd}}$ & Year & $N$ & $\Gamma$ & $G(\bpi)$ & $\pi_{\text{hyd}}$ \\
\midrule
1988 & 10 & 0.563 & 0.442 & 0.025 & 2004 & 10 & 0.524 & 0.385 & 0.057 \\
1989 & 10 & 0.565 & 0.429 & 0.029 & 2005 & 10 & 0.502 & 0.366 & 0.074 \\
1990 & 10 & 0.567 & 0.420 & 0.037 & 2006 & 11 & 0.518 & 0.424 & 0.064 \\
1991 & 10 & 0.578 & 0.431 & 0.030 & 2007 & 10 & 0.489 & 0.350 & 0.087 \\
1992 & 10 & 0.571 & 0.417 & 0.030 & 2008 & 10 & 0.486 & 0.341 & 0.081 \\
1993 & 10 & 0.573 & 0.419 & 0.026 & 2009 & 10 & 0.504 & 0.353 & 0.077 \\
1994 & 10 & 0.567 & 0.413 & 0.022 & 2011 & 11 & 0.478 & 0.369 & 0.088 \\
1995 & 10 & 0.578 & 0.411 & 0.025 & 2012 & 11 & 0.506 & 0.379 & 0.079 \\
1996 & 10 & 0.566 & 0.407 & 0.026 & 2013 & 11 & 0.484 & 0.371 & 0.085 \\
1997 & 10 & 0.553 & 0.409 & 0.027 & 2014 & 11 & 0.361 & 0.352 & 0.117 \\
\cmidrule{6-10}
1998 & 10 & 0.541 & 0.449 & 0.017 & 2017 & 11 & 0.506 & 0.446 & 0.045 \\
1999 & 10 & 0.642 & 0.425 & 0.035 & 2018 & 11 & 0.493 & 0.444 & 0.042 \\
2000 & 10 & 0.494 & 0.420 & 0.036 & 2019 & 11 & 0.467 & 0.445 & 0.039 \\
2001 & 10 & 0.493 & 0.410 & 0.038 & 2020 & 11 & 0.418 & 0.463 & 0.035 \\
2002 & 11 & 0.501 & 0.447 & 0.039 & 2022 & 11 & 0.400 & 0.475 & 0.026 \\
2003 & 11 & 0.514 & 0.444 & 0.047 & 2024 & 11 & 0.419 & 0.470 & 0.021 \\
\bottomrule
\end{tabular}
\end{center}
(2021: $0.432$, $0.462$, $0.028$; 2023: $0.381$, $0.470$, $0.030$. The rule separates base 1990 from base 2017.)

\textbf{Outliers.} On base 1990 the $z$-scores of $\Gamma$ exceed $1.5$ in absolute value in two years: 1999 ($\Gamma = 0.642$, $z = 2.15$) and 2014 ($0.361$, $z = -3.14$). In 2014 the slow mode $\mathbf{v}_2$ is concentrated on hydrocarbons ($1.0$, against at most $0.45$ in absolute value elsewhere). The hydrocarbons self-purchase share jumps from $0.478$ (2013) to $0.658$ (2014), the largest diagonal change of any sector. A near-absorbing diagonal entry creates an eigenvalue near the diagonal value, so the 2014 $\Gamma$ reflects one cell of the table, not a disintegration of the economy. It is most likely a compilation artefact of the last base-1990 table.

\textbf{Why not compare across the change of base.} The 2017 base has a new classification (42 products, 39 industries against 35), a new SNA vintage and new valuation conventions. $\Gamma$, $G(\bpi)$ and $\pi_{\text{hyd}}$ depend on aggregation and conventions (Exercise~17.1), so a jump at the break mixes measurement with structure. A further caveat: on base 1990, $N$ alternates between 10 and 11, because public administration has no intermediate sales in diagonal-make years. Changes between, say, 2005 ($N = 10$) and 2006 ($N = 11$) are partly a change of node set.

\solpart{b} \textbf{$\sigma_B$ by year: not computable here} (EH microdata not fetched). Method: Exercise~17.2 applied to EH 2000, 2005, 2011 (there is no 2010 survey), 2015 and 2019, with a fixed partition and fixed definitions of formality; then plot. INE's published Gini fell from $0.600$ (2005) to $0.426$ (2019) (Chapter~16), so the total variance fell. Whether $\sigma_B$ fell too depends on whether the fall was between or within blocks; decompositions attributing it to labour income at the bottom \citep{vargas2015explaining} suggest a fall in $V_B$.

\solpart{c} \textbf{Structural breaks in $\Gamma$.} Chow tests on the annual base-1990 series, 1988--2013 (2014 excluded), with a model of level plus linear trend and a break in both:
\begin{center}
\begin{tabular}{lcccc}
\toprule
& \multicolumn{2}{c}{All years ($n = 25$)} & \multicolumn{2}{c}{Excluding 1999 ($n = 24$)} \\
Break at & $F(2, n-4)$ & $p$ & $F$ & $p$ \\
\midrule
2006 (nationalization) & 0.108 & 0.898 & 0.940 & 0.407 \\
2009 (constitution) & 0.170 & 0.845 & 1.193 & 0.324 \\
\bottomrule
\end{tabular}
\end{center}
No break is detected at either date, and the OLS trend is $-0.042$ per decade ($-0.041$ without 1999). The $p$-values come from the $F$ distribution, integrated numerically; the integration is checked against the 5\% critical value of $F(2,20)$, which is $3.49$. Consistent with Example~\ref{ch17:ex:trajectory}: the nationalization changed revenue flows, not $W^P$. $\pi_{\text{hyd}}$ shows the gas boom (from $0.047$ in 2003 to $0.088$ in 2011), but that rise starts before 2006.

\solpart{d} \textbf{Projection of $\sigma_B$.} No data. With the book's \emph{illustrative} trajectory, $0.68$ (1990) to $0.62$ (2019), the trend is $-0.021$ per decade, and $\sigma_B$ would reach $0.40$ around 2125. That number illustrates the arithmetic only; a projection needs the EH series of part~(b) and a model of why $\sigma_B$ moves.

%% ---------------------------------------------------------------------
\exercisesolution{17.6}{Dashboard comparison}

\solpart{a} \textbf{Measured instruments} (2019, book12\_isic, domestic, input shares; INE for Bolivia, OECD ICIO 2023 for Chile and Costa Rica):
\begin{center}
\begin{tabular}{lccccc}
\toprule
& $\Gamma$ & $G(\bpi)$ & $\Neff(W^P)$ of 11 & ECI 2019 & Gini (WDI) \\
\midrule
Bolivia & 0.475 & 0.435 & 6.67 & $-1.295$ & 41.6 (2019) \\
Chile & 0.557 & 0.568 & 4.22 & $-0.059$ & 45.3 (2017) \\
Costa Rica & 0.469 & 0.611 & 4.14 & $0.257$ & not in \nolinkurl{data/processed} \\
\bottomrule
\end{tabular}
\end{center}
Costa Rica's ECI is from the Growth Lab rankings file in \nolinkurl{data/raw/atlas}. The largest power weights are other services in Chile ($0.429$) and in Costa Rica ($0.414$). $\sigma_B$ needs household microdata (CASEN for Chile, ENAHO for Costa Rica, EH for Bolivia); not fetched.

\solpart{b} $\Neff(W^I)$: not computable (audience data not fetched). Method of Exercise~17.4 on a common four-ecosystem partition: state or public, private commercial, community, social media.

\solpart{c} \textbf{Dashboard and ``healthiest''.} Rescaled bars (benchmark $= 1$):
\begin{itemize}
\item $\Gamma$: Bolivia $0.475$, Chile $0.557$, Costa Rica $0.469$;
\item $1 - G(\bpi)$: $0.565$, $0.432$, $0.389$;
\item $\Neff/N$: $0.606$, $0.384$, $0.377$.
\end{itemize}
Chile is the most integrated, Bolivia the least concentrated. On both production-network coordinates Costa Rica, the most complex of the three, is the furthest from the benchmark. ``Healthiest'' is therefore undefined without weights on independent coordinates (\nolinkurl{docs/conventions.md}, \S4), and the measured instruments rank the countries opposite to complexity. As in Chapter~17, a concentrated $\bpi(W^P)$ marks a service-hub structure, not exclusion.

\solpart{d} \textbf{Recommendation.} The largest deficit depends on the coordinate:
\begin{itemize}
\item Bolivia: complexity (ECI $-1.30$). The recommended operators are $\mathcal{L}_{\text{info}}$ and $\mathcal{L}_{\text{HC}}$ aimed at related products, not $\Gamma$-raising operators.
\item Costa Rica, on power concentration: an operator that diversifies suppliers away from the services hub. Proposition~\ref{ch15:prop:political_cost} says integration is free politically, but redistribution is not.
\end{itemize}
Choose the operator by the instrument, then sequence by Conjecture~\ref{ch15:thm:sequencing}.

%% ---------------------------------------------------------------------
\exercisesolution{17.7}{Bolivia's escape plan}

\solpart{a} \textbf{Current position.} $\Gamma = 0.475$ (like-for-like) and $G(\bpi) = 0.435$, both measured. ECI $-1.30$ in 2019, down from $-1.21$ in 2006. $\sigma_B$, $\Neff(W^I)$ and $\Psi$ remain illustrative ($0.62$, $3.4$, $1.56$).

\solpart{b} \textbf{Binding constraint.} Among the \emph{measured} instruments, none is far from comparators: $\Gamma$ sits inside the range, and $G(\bpi)$ is the lowest of 23 economies. The binding constraint is complexity, which is not a dashboard instrument but is the memo row of Example~\ref{ch17:ex:dashboard}. Among the illustrative rows, $\sigma_B$ and $\Neff(W^I)$ would be the candidates if the microdata confirmed them. Because $\Gamma$ and $G(\bpi)$ are independent coordinates, an operator chosen to raise $\Gamma$ does not address either.

\solpart{c} \textbf{Sequence:} information, then institutions, education, infrastructure and ISI-B (Conjecture~\ref{ch15:thm:sequencing}), aimed at products related to the existing basket (Exercise~16.7(b)). Exercise~15.4 shows when this order beats the reverse: only if later operators are complementary to earlier ones.

\solpart{d} \textbf{Political feasibility.} Integration and complexity operators need not cost any coalition power (Proposition~\ref{ch15:prop:political_cost}), so a coalition of the middle class and workers ($\{M, W\}$), or of the bourgeoisie, middle class and workers ($\{B, M, W\}$), can support them without internal conflict. The redistributive part, shifting rents or power from the hydrocarbons--state nexus, is where the coalition's own members' $\Delta\pi_i \ge 0$ constraints bind (Exercise~15.7). A $\{B, M, W\}$ coalition is compatible if the plan's surplus share $\Delta\Omega/N$ compensates $B$ for any power loss (Proposition~\ref{ch14:prop:power_gain}).

\solpart{e} \textbf{Simulation.} Use the three-block model of Example~\ref{ch10:ex:bolivia_three} with $\rho = 0.85$ and $\mu = 0.2$. The script reproduces its steady state ($\theta^* = 1.1245$, $1.0411$, $0.9871$; $\rho(J) = 0.866$). Apply an information operator: cross-block $W^I$ weight per pair from $0.025$ to $0.05$, with own-block weights cut equally.
\begin{itemize}
\item \emph{Long run}: $\theta^* = (1.1009,\allowbreak 1.0439,\allowbreak 1.0062)$; the gas--rural preference gap falls from $0.1374$ to $0.0947$. The variance of allocations falls from $0.03890$ to $0.03229$; its power part ($0.01985$) is untouched, because $\bpi$ is.
\item $\Gamma(W^I)$ rises from $0.225$ to $0.450$. $\Gamma(W^P) = 0.090$ and $\Neff(\bpi) = 8.34$ do not move.
\item \emph{30 periods, generational} ($\rho = 0.85$; one period is one generation, Chapter~10): $14.0\%$ of the adjustment occurs after 1 period, $77.9\%$ after 10, $98.9\%$ after 30.
\item \emph{30 years, annual} ($\rho_a = 0.85^{1/25}$, with $\mu_a = 0.00864$ chosen to keep $\mathcal{M}$ fixed; $\rho(J) = 0.9940$): only $0.6\%$ after one year, $5.9\%$ after ten and $16.6\%$ after thirty. The gap falls from $0.1374$ to $0.1303$.
\end{itemize}
\textbf{Answer.} Target the information structure ($\Neff(W^I)$, the cross-block weight of $W^I$) and complexity, not $\Gamma(W^P)$, which is already in the comparator range. With generational persistence $\rho = 0.85$, thirty \emph{years} buy about a sixth of the long-run preference convergence; most of the effect takes several generations. Plans should be judged against that horizon.

\emph{Common mistake:} simulating 30 periods of the generational model and calling them years, which overstates the speed of formation by a factor of about six here ($98.9\%$ against $16.6\%$).

\endgroup
\end{document}
